% Generated draft: requires mathematical and source review.
\documentclass[11pt]{article}
\usepackage[margin=25mm]{geometry}
\usepackage{fontspec}
\setmainfont{Latin Modern Roman}
\usepackage{amsmath,amssymb,mathtools,unicode-math}
\setmathfont{Latin Modern Math}
\usepackage{enumitem,needspace,xurl,hyperref}
\setlength{\parindent}{0pt}
\setlength{\parskip}{5pt}
\setlength{\emergencystretch}{3em}
\newcommand{\N}{\mathbb{N}}
\newcommand{\Z}{\mathbb{Z}}
\newcommand{\Q}{\mathbb{Q}}
\newcommand{\R}{\mathbb{R}}
\newcommand{\C}{\mathbb{C}}
\newcommand{\F}{\mathbb{F}}
\newcommand{\A}{\mathbb{A}}
\newcommand{\PP}{\mathbb{P}}
\newcommand{\E}{\mathbb{E}}
\newcommand{\eps}{\varepsilon}
\newcommand{\dd}{\,\mathrm{d}}
\begin{document}
\section*{UnsolvedMath ID: 30004667}
\textbf{Source problem code:} catalogue entry
\subsection{Definitions and mathematical statement}
The supplied record asks about an equivariant-multiplicity formula involving a Duistermaat--Heckman measure. It does not state that formula or specify the geometric objects, acting group, hypotheses, normalization conventions, or class of examples. Thus neither the relevant equivariant multiplicity nor the measure can be defined in its intended setting from the supplied information alone.

The two requested tasks are:
\begin{enumerate}
\item Prove the equivariant-multiplicity formula referred to in the source report without using Duistermaat--Heckman measures.
\item Compute concrete examples of the equivariant multiplicities occurring in that formula.
\end{enumerate}
The formula, its hypotheses, and the intended examples remain unspecified. These tasks therefore record the source's intended questions but do not constitute a fully specified mathematical assertion. In particular, the permitted methods in a proof ``without using'' these measures are not defined; this condition does not by itself require an elementary proof.

According to the supplied dated curation, the report leaves both questions unresolved, describes known computation methods as not completely algorithmic, and gives only a rank-one example. The example and its ambient setting are not supplied here. These descriptions are source qualifications, not independently verified mathematical conclusions.
\subsection{Short English statement}
The source asks for a proof of its equivariant-multiplicity formula avoiding Duistermaat--Heckman measures, and for concrete computations. The actual formula, hypotheses, and example class are missing, so a precise mathematical formulation requires further source information. The supplied open status requires independent verification.
\subsection{Sources}
\begin{itemize}
\item Supplied local registry record: permanent UnsolvedMath ID \(30004667\), source problem code \texttt{catalogue entry}, title \emph{Elementary Proofs and Computations of Equivariant Multiplicities}, in the Oberwolfach Reports Open Problems collection. The record gives proposed year 2021 and names no proposer.
\item Supplied report citation: \emph{Algebraic Groups (2022)}, Oberwolfach Reports, pp.~1087--1148, DOI \texttt{10.4171/owr/2021/20}. \url{https://doi.org/10.4171/owr/2021/20}. The supplied title year 2022 and report designation 20/2021 are retained without resolving their relationship.
\item Supplied literature source: \emph{Oberwolfach Report 20/2021}. \url{https://ems.press/content/serial-article-files/46898}.
\item Supplied literature source: P\u{a}durariu, \emph{K-theoretic Hall algebras}. \url{https://arxiv.org/abs/1911.05526}. Its relationship to the unspecified formula is not explained in the record.
\item Literature assessment dated exactly \texttt{2026-08-22}: the supplied curation labels the problem open and says that both questions remain unresolved in the report, known computation methods are not completely algorithmic, and only a rank-one example is given, with no general solution established. This is supplied dated curation, not an independently verified assessment; the sources and time-sensitive claims have not been checked here.
\end{itemize}
\end{document}
